\documentclass[twocolumn]{article}

% DAFx-style packages
\usepackage[utf8]{inputenc}
\usepackage[T1]{fontenc}
\usepackage{times}
\usepackage{amsmath,amssymb}
\usepackage{graphicx}
\usepackage{url}
\usepackage{hyperref}
\usepackage{booktabs}
\usepackage{xcolor}
\usepackage{listings}
\usepackage[a4paper,left=19mm,right=19mm,top=25.4mm,bottom=25.4mm]{geometry}

% Code listing style
\lstset{
  basicstyle=\ttfamily\footnotesize,
  breaklines=true,
  frame=single,
  xleftmargin=2em,
  framexleftmargin=1.5em
}

% Title formatting
\makeatletter
\renewcommand{\maketitle}{%
  \twocolumn[%
    \begin{center}
      {\Large\bfseries\MakeUppercase{\@title}\par}
      \vspace{1em}
      {\normalsize\@author\par}
      \vspace{1.5em}
    \end{center}
  ]%
}
\makeatother

\title{Real-Time State-Space Parameterization in Digital Equalization:\\A Production-Grade SVF Implementation}

\author{%
  \textbf{Gary Doman}\\
  FreeEQ8 / ProEQ8 Open-Source DSP Project\\
  \url{https://github.com/GareBear99/FreeEQ8}
}

\date{}

\begin{document}

\maketitle

\begin{abstract}
We present a production-grade parametric equalizer implementing the Simper State Variable Filter (SVF) topology to eliminate high-frequency magnitude cramping without oversampling. The system achieves exact cutoff placement at 16\,kHz (where traditional RBJ biquads exhibit 199\% Q distortion) while consuming only 0.62\% of a single-core CPU budget at 44.1\,kHz. A lock-free SPSC triple-buffer isolates the audio thread from UI rendering, and a variable-cadence coefficient engine reduces Dynamic EQ overhead by 80\%. The implementation is available as FreeEQ8 (8-band, GPL-3.0) and ProEQ8 (24-band, commercial) from a single codebase. Live demonstration accompanies this paper.
\end{abstract}

\section{Introduction}

Traditional parametric EQ implementations use Robert Bristow-Johnson's Audio EQ Cookbook biquad formulas~\cite{rbj} in Transposed Direct Form II (TDF-II). While mathematically correct, the bilinear transform introduces frequency cramping near Nyquist: a Bell filter at 16\,kHz with $Q=1.0$ at 44.1\,kHz exhibits an effective bandwidth \textbf{199\% narrower} than intended.

Industry solutions include brute-force oversampling (adding latency and CPU cost) or proprietary polynomial analog-matching curves. This paper documents a third path: the Simper SVF topology~\cite{simper}, which pre-warps the cutoff frequency via $g = \tan(\pi \cdot f_c / f_s)$, achieving exact cutoff placement at any frequency up to Nyquist without oversampling.

\section{Filter Topology}

\subsection{RBJ TDF-II (FreeEQ8)}

The transfer function follows the standard biquad form:
\begin{equation}
H(z) = \frac{b_0 + b_1 z^{-1} + b_2 z^{-2}}{1 + a_1 z^{-1} + a_2 z^{-2}}
\end{equation}

Measured Q distortion at 44.1\,kHz sample rate is shown in Table~\ref{tab:qdistortion}.

\begin{table}[h]
\centering
\caption{RBJ Q distortion (Bell, $Q=1.0$, 44.1\,kHz)}
\label{tab:qdistortion}
\begin{tabular}{lcc}
\toprule
Frequency & Effective Q & Error \\
\midrule
1\,kHz & 1.005 & +0.5\% \\
8\,kHz & 1.337 & +33.7\% \\
16\,kHz & 2.990 & +199\% \\
\bottomrule
\end{tabular}
\end{table}

\subsection{Simper SVF (ProEQ8)}

The SVF uses pre-warped coefficients~\cite{simper}:
\begin{align}
g &= \tan(\pi \cdot f_c / f_s) \\
k &= 1/Q \\
a_1 &= 1 / (1 + g \cdot (g + k))
\end{align}

Per-sample processing follows the optimized bounded form:
\begin{lstlisting}[language=C++]
v3 = v0 - ic2eq;
v1 = a1*ic1eq + a2*v3;  // bandpass
v2 = ic2eq + a2*ic1eq + a3*v3; // lowpass
ic1eq = v1 + v1 - ic1eq;
ic2eq = v2 + v2 - ic2eq;
out = m0*v0 + m1*v1 + m2*v2;
\end{lstlisting}

All eight filter types (Bell, LowShelf, HighShelf, LP, HP, Bandpass, Notch, AllPass) emerge from a single two-integrator core with different output mix coefficients $(m_0, m_1, m_2)$.

\subsection{Measured Magnitude Response}

To validate de-cramping, we swept a sine (20\,Hz--20\,kHz) through Bell filters (+6\,dB, $Q=1.0$) at $f_c = 16$\,kHz. Results in Table~\ref{tab:magnitude}.

\begin{table}[h]
\centering
\caption{Magnitude at $f_c = 16$\,kHz, 44.1\,kHz sample rate}
\label{tab:magnitude}
\begin{tabular}{lcc}
\toprule
Topology & Magnitude & Error \\
\midrule
RBJ @ 44.1\,kHz & +0.73\,dB & $-5.27$\,dB \\
SVF @ 44.1\,kHz & +6.00\,dB & 0.00\,dB \\
RBJ @ 4$\times$ OS & +4.82\,dB & $-1.18$\,dB \\
\bottomrule
\end{tabular}
\end{table}

The SVF achieves exact gain without oversampling. RBJ@4$\times$OS reduces error but costs 4$\times$ CPU.

\section{Real-Time Safety Architecture}

\subsection{SPSC Triple-Buffer}

Three buffer slots indexed by \texttt{\{writeSlot, midSlot, readSlot\}}. Audio thread writes into \texttt{writeSlot}; atomically swaps with \texttt{midSlot} using \texttt{memory\_order\_release}. UI thread reads by swapping \texttt{midSlot} $\leftrightarrow$ \texttt{readSlot} with \texttt{memory\_order\_acquire}. Zero mutex, zero blocking.

Stress testing on Intel Ivy Bridge (2012): \textbf{0 data tears across 239M samples}.

\subsection{Variable-Cadence Dynamic EQ}

Dynamic EQ recomputes coefficients per-sample. We observe that during sustained notes, envelope changes of $<0.1$\,dB per sample are inaudible within a 4-sample batch (0.09\,ms at 44.1\,kHz).

\begin{lstlisting}
delta = |dynGainMod - lastDynGainMod|
if delta > 0.1 dB:
    update coefficients (per-sample)
else if intervalCounter++ >= 4:
    update coefficients (batched)
\end{lstlisting}

Measured savings: \textbf{80\% reduction} in coefficient updates with no audible difference. Transients immediately restore per-sample accuracy.

\section{Benchmarks}

All benchmarks from standalone C++17, \texttt{g++ -O3}, 44.1\,kHz, 512-sample blocks.

\begin{table}[h]
\centering
\caption{Single-instance filter cost (8-band stereo)}
\label{tab:benchmarks}
\begin{tabular}{lccc}
\toprule
Path & ns/sample & CPU\% & Headroom \\
\midrule
RBJ 8-band & 41.0 & 0.36\% & 277$\times$ \\
SVF 8-band & 72.7 & 0.63\% & 161$\times$ \\
SVF DynEQ & 68.8 & 0.61\% & 165$\times$ \\
\bottomrule
\end{tabular}
\end{table}

Instance scaling (1$\rightarrow$128 instances): per-instance cost rises only 5\%. At 128 SVF instances, total CPU = 85.8\% of one core. A modern 8-core CPU can host $\sim$900 instances.

\section{Product Architecture}

The codebase produces two plugins via \texttt{\#if PROEQ8}:

\textbf{FreeEQ8} (GPL-3.0): 8-band RBJ biquad. Zero real-time restrictions. Offline export limited to 4:30.

\textbf{ProEQ8} (\$20): 24-band SVF with de-cramped response. Adds saturation modes, A/B comparison, auto-gain, collision detection. No export limit.

Both validated via \texttt{pluginval} at strictness-level-10 on macOS/Linux/Windows CI.

\section{Demonstration}

The live demonstration showcases:
\begin{enumerate}
\item Side-by-side RBJ vs SVF response at 16\,kHz
\item Real-time spectrum analyzer with lock-free rendering
\item Dynamic EQ envelope tracking on transient material
\item 24-band surgical EQ workflow in ProEQ8
\end{enumerate}

Source code, benchmarks, and builds available at:\\
\url{https://github.com/GareBear99/FreeEQ8}

\section{Conclusion}

We have demonstrated that the Simper SVF topology eliminates high-frequency cramping while maintaining sub-1\% CPU overhead. The lock-free architecture ensures glitch-free operation across all tested DAWs, and variable-cadence Dynamic EQ reduces computational cost by 80\% without perceptual degradation.

\begin{thebibliography}{9}

\bibitem{simper}
A. Simper, ``Solving the continuous SVF equations using trapezoidal integration and equivalent currents,'' Cytomic, 2013.
\url{https://cytomic.com/files/dsp/SvfLinearTrapOptimised2.pdf}

\bibitem{rbj}
R. Bristow-Johnson, ``Audio EQ Cookbook,'' musicdsp.org, 1994.
\url{https://www.musicdsp.org/files/Audio-EQ-Cookbook.txt}

\bibitem{pirkle}
W. Pirkle, \textit{Designing Audio Effect Plugins in C++}, Focal Press, 2019.

\bibitem{reiss}
J. Reiss and A. McPherson, \textit{Audio Effects: Theory, Implementation and Application}, CRC Press, 2014.

\end{thebibliography}

\end{document}
